\documentclass[11pt,a4paper]{article}
\usepackage[margin=1in]{geometry}
\usepackage{amsmath,amssymb,amsfonts}
\usepackage{mathtools}
\usepackage{lmodern}
\usepackage[T1]{fontenc}
\usepackage[utf8]{inputenc}
\usepackage{textcomp}
\usepackage{microtype}
\usepackage{parskip}
\usepackage{enumitem}
\usepackage{array}
\usepackage{booktabs}
\usepackage{hyperref}
\hypersetup{hidelinks,
  pdftitle={From Primitives to Metamaterials and Photonics},
  pdfauthor={Allen Proxmire}}
\setlength{\parindent}{0pt}
\setlength{\parskip}{6pt plus 2pt minus 1pt}

\title{\vspace{-1cm}\Large\bfseries From Primitives to Metamaterials and Photonics\\[6pt]
\large\mdseries A Walkthrough of the Event Density Derivation}
\author{Allen Proxmire}
\date{May 2026}

\begin{document}
\maketitle

\section{The Question}

\subsection{What this walkthrough derives}

This walkthrough derives, from substrate primitives, the structural backbone of four Nobel-frontier results in metamaterials and photonics:

\begin{enumerate}[noitemsep]
\item \textbf{The effective constitutive tensors} $\varepsilon_{\rm eff}^{ij}(\omega)$ and $\mu_{\rm eff}^{ij}(\omega)$ of a periodic subwavelength microstructure, derived via a two-scale homogenization expansion and a cell problem on a unit cell $Y$. The cell-problem solution defines a corrector field $\chi^j(y)$ whose cell average provides the substrate-level correction to the naive arithmetic average of the microstructure response.
\item \textbf{The negative refractive index} $n_{\rm eff}(\omega) = -\sqrt{\varepsilon_r(\omega)\mu_r(\omega)}$ when both effective tensors are negative in an overlapping frequency band. The negative branch is FORCED by causal outgoing-wave selection (Pendry 2000). The Drude form $\varepsilon_{\rm eff}(\omega) = \varepsilon_\infty(1 - \omega_p^2/\omega^2)$ from wire-array microstructures and the Lorentz form $\mu_{\rm eff}(\omega) = 1 - F\omega^2/(\omega^2 - \omega_0^2 + i\gamma\omega)$ from split-ring resonators are each derived as cell-problem outputs.
\item \textbf{The constitutive transformation under substrate deformation} $\varepsilon'^{ij} = (\det J)^{-1} J^i{}_k J^j{}_l \varepsilon^{kl}$ and the coarse-grained \textbf{effective metric} $g_{ij} = (J^{-1})^k{}_i (J^{-1})^l{}_j \delta_{kl}$ on the wave equation, with the wave equation rewritten in covariant Laplace-Beltrami form. The metric is a coarse-grained rule-type metric on the substrate, not spacetime curvature.
\item \textbf{The Pendry spherical cloak} (Pendry 2006). The radial map $r' = R_1 + \alpha r$ with $\alpha = (R_2 - R_1)/R_2$ produces a shell-supported anisotropic medium whose effective metric vanishes radially at the inner shell. The cloaked interior is topologically excluded from the substrate seen by the propagating channel.
\item \textbf{The metasurface generalized Snell's law} (Capasso 2011). An engineered codimension-1 surface carrying a tangential phase profile $\Phi(r_\parallel)$ imposes a tangential momentum kick $k_{\parallel,2} = k_{\parallel,1} + \nabla_\parallel\Phi$, yielding
\[
n_1 \sin\theta_i = n_2 \sin\theta_t + \frac{\lambda}{2\pi}\,\frac{d\Phi}{dx}.
\]
\item \textbf{The photonic bandgap} (Yablonovitch 1987) as a frequency window in which no Bloch eigenmode exists on the periodic substrate at $a \sim \lambda$. The relevant Bloch theorem is derived inline in compact form.
\item \textbf{The unifying structural identity}:
\[
\text{Metamaterials} = \text{engineered microstructure} + \text{engineered deformation} + \text{engineered discontinuity}.
\]
The three operations exhaust the substrate-level wave-control space within the coarse-graining window $\ell_P \ll a \ll \lambda \ll L$.
\end{enumerate}

Each result is derived, not posited. Every required two-scale expansion step, cell-problem statement, Jacobian computation, effective-metric identification, jump condition, and Bloch-eigenmode statement appears inline.

\subsection{Why standard photonics treats $\varepsilon$, $\mu$, and metamaterial response phenomenologically}

Conventional macroscopic electromagnetism treats $\varepsilon$, $\mu$, and the refractive index $n$ as phenomenological material parameters fit to experiment. The Maxwell equations are taken as exact at the macroscopic level, with material response encoded in constitutive relations $\mathbf{D} = \varepsilon\mathbf{E}$ and $\mathbf{B} = \mu\mathbf{H}$. The standard derivations of $\varepsilon$ and $\mu$ from molecular polarizabilities (Clausius-Mossotti, Lorentz local-field corrections) treat the microscopic constituents as fixed atoms or molecules whose responses are summed over a coarse-graining volume.

The recognition that subwavelength structuring can engineer $\varepsilon_{\rm eff}$ and $\mu_{\rm eff}$ unattainable in natural materials --- including negative values --- was Pendry's central insight. The further recognition that smooth coordinate transformations of Maxwell's equations correspond to anisotropic constitutive tensors that can be physically built by metamaterials was the transformation-optics insight. The realization that engineered codimension-1 phase profiles produce anomalous refraction outside the bulk Snell relation was the metasurface insight. The recognition that periodic dielectric structures at $a \sim \lambda$ produce band structures with frequency gaps was the photonic-bandgap insight.

Standard theory delivers the macroscopic equations correctly. It is silent about what $\varepsilon$ and $\mu$ are ontologically: it treats them as numbers attached to a material, with no account of \emph{why} a particular microstructure produces a particular response, \emph{why} cloaking works mechanistically rather than merely formally as a coordinate identity, \emph{what} a metasurface boundary condition is at the level of substrate ontology, or \emph{why} the four Nobel-frontier directions share a common structural skeleton.

\subsection{What Event Density claims}

In the ED framework the substrate carries rule-type structure at scale $\ell_P$, where rule-type means the discrete pattern of participation-rule alignment that the substrate locally enforces on channels of ED-flow. Coarse-graining within the window $\ell_P \ll a \ll \lambda \ll L$ produces an effective wave equation whose coefficients are coarse-grained projections of the substrate's local rule-type content. Three independent engineered modulations of this content control wave propagation:

\begin{itemize}[noitemsep]
\item \textbf{Smooth periodic microstructure} at scale $a \ll \lambda$ produces effective constitutive tensors via homogenization. The substrate is unstrained; only its internal rule-type content is modulated. Pendry-negative-index lives here.
\item \textbf{Smooth deformation} $u(X)$ produces an effective metric and transformed constitutive tensors via transformation optics. The substrate is strained; its rule-type content is transported covariantly along the deformation. Pendry-cloak lives here.
\item \textbf{Engineered codimension-1 discontinuity} with engineered phase profile produces tangential momentum kicks at the interface. The substrate is unstrained in the bulk; only the interface carries engineered rule-type content. Capasso-metasurface lives here.
\item \textbf{Periodic microstructure at $a \sim \lambda$}, adjacent to the homogenization regime, supports Bloch eigenmodes with band structure. Yablonovitch-bandgap lives here.
\end{itemize}

The substrate-level statement, in one sentence: \textbf{metamaterials are engineered modulations of the substrate's coarse-grained rule-type structure; the three independent modulation axes (microstructure, deformation, discontinuity) exhaust the wave-control space within the coarse-graining window.}

\subsection{The chain in summary}

The derivation chain runs: substrate primitives P-MM-1 through P-MM-6 (\S2) $\to$ two-scale lift, derivative splitting, order-by-order equations, cell problem, corrector field, averaging operator, effective-tensor formula (\S3) $\to$ identification with $\varepsilon$ and $\mu$, Drude wire-array response, Lorentz split-ring response, doubly-negative band, FORCED negative-index branch (\S4) $\to$ deformation Jacobian, rule-type deformation tensor, constitutive transformation, effective metric in covariant form, Pendry spherical cloak, topological exclusion of cloaked interior (\S5) $\to$ codimension-1 discontinuity, phase-imprinting boundary condition, tangential momentum kick, generalized Snell's law (\S6) $\to$ Bloch form, band structure, bandgap formation, relation to homogenization (\S7) $\to$ unifying substrate-level statement, three-axis exhaustiveness, scope limits (\S8) $\to$ FORCED / INHERITED / OPEN accounting (\S9) $\to$ exact claims (\S10).

\section{The Substrate Primitives}

Six substrate primitives are used. No new primitives are introduced anywhere in the derivation; the entire walkthrough runs on this list.

\subsection{P-MM-1. Substrate rule-type microstructure}

The substrate carries rule-type structure at scale $\ell_P$ --- discrete participation rules indexed by spatial position and internal type, encoding the substrate's local commitment patterns for ED-flow. At coarse-grained scales $a \gg \ell_P$, engineered inclusions modulate this rule-type structure, producing a position-dependent coarse-grained response. The coefficients $\varepsilon^{ij}(x)$, $(\mu^{-1})^{ij}(x)$, and $n^2(x)$ of the macroscopic wave equation are coarse-grained projections of this rule-type content.

\textbf{Substrate-level meaning.} $\varepsilon$ and $\mu$ are not fundamental fields. They are coarse-grained labels for the substrate's local rule-type response: $\varepsilon^{ij}(x)$ is the coarse-grained polarizability of the rule-type structure at $x$ --- the local tendency to align ED-flow under an applied electric tension --- and $(\mu^{-1})^{ij}(x)$ is the coarse-grained inverse circulation susceptibility --- the local resistance to closed-loop ED-flow under an applied magnetic tension.

\textbf{Engineered modulation.} Replacing the substrate's vacuum rule-type content with an engineered pattern of inclusions modulates these coefficients. The pattern of inclusions, coarse-grained, is the metamaterial response.

\subsection{P-MM-2. Subwavelength periodicity}

When inclusions are arranged with period $a$ satisfying $a \ll \lambda$, the rule-type structure admits a two-scale description in a slow coordinate $X = x$ and a fast coordinate $y = x/a$, with $y$ taking values on the unit cell $Y = [0, 1]^3$ (in normalized fast coordinates). Field quantities lift to two-scale form $\psi(x) \to \tilde\psi(X, y)$ with $\tilde\psi$ periodic in $y$. The two-scale expansion is the substrate-level mathematical reflection of the scale separation $a/\lambda \ll 1$.

\textbf{Periodic substrate.} The rule-type content is the same on each unit cell up to translations by $aR$ with $R$ in the integer lattice $\mathbb{Z}^3$. Practical metamaterials approximate this idealization to within a few percent at microwave frequencies.

\subsection{P-MM-3. ED-gradient deformation}

A slowly varying displacement field $u(X)$ acting on the substrate produces a coarse-grained deformation
\[
x'^i = x^i + u^i(X),
\]
whose Jacobian
\[
J^i{}_j = \delta^i{}_j + \frac{\partial u^i}{\partial X^j}
\]
transports rule-type structure covariantly. The deformation is smooth and orientation-preserving: $\det J > 0$ everywhere, $J$ and $J^{-1}$ bounded. Time-independent deformations carry no internal dispersion of their own; all dispersion enters through the constitutive coefficients at the microstructure level.

\textbf{Substrate-level meaning.} The deformation is a coordinate-level engineered displacement of the substrate. Rule-type tensors transform as densities of weight one: a coefficient $\varepsilon^{ij}(x)$ at the undeformed point $x$ becomes the coefficient $\varepsilon'^{ij}(x')$ at the deformed point $x'$, with the transformation rule derived in \S5.

\subsection{P-MM-4. Channel propagation in structured media}

Wave propagation in the coarse-grained substrate is governed by an effective second-order wave equation
\[
\partial_{x^i}\!\left[A^{ij}(x)\,\partial_{x^j}\psi(x)\right] + k_0^2\,n^2(x)\,\psi(x) = 0,
\]
where $\psi$ is a scalar component of the electromagnetic field (or any field reducible to scalar form for a chosen polarization), $A^{ij}(x)$ is a coarse-grained rank-2 tensor coefficient determined by the substrate's local rule-type content, $n^2(x)$ is a coarse-grained scalar coefficient, and $k_0 = \omega/c$ is the vacuum wavenumber.

\textbf{Polarization identifications.} For Maxwell's equations in the time-harmonic regime, the scalar form corresponds to:
\begin{itemize}[noitemsep]
\item TM modes (transverse magnetic, $\mathbf{E}$ in the plane of incidence): $A^{ij} \leftrightarrow (\mu^{-1})^{ij}$, $n^2 \leftrightarrow \varepsilon$.
\item TE modes (transverse electric, $\mathbf{H}$ in the plane of incidence): $A^{ij} \leftrightarrow \varepsilon^{ij}$, $n^2 \leftrightarrow \mu$.
\end{itemize}

The derivations of \S3 apply to either polarization by reading the identifications.

\subsection{P-MM-5. Interface rule-type discontinuity}

A codimension-1 surface $\Sigma$ across which the rule-type microstructure or deformation jumps imposes jump conditions on the coarse-grained wave field. In the absence of engineered surface phase profile or surface currents, the standard tangential-field continuity holds: $\mathbf{E}_\parallel$ and $\mathbf{H}_\parallel$ are continuous across $\Sigma$. An engineered surface phase profile $\Phi(r_\parallel)$ imposes a phase-jump matching condition derived in \S6.

\textbf{Substrate-level meaning.} A metasurface is a codimension-1 engineered rule-type discontinuity: the substrate carries different rule-type content on the two sides of $\Sigma$, with the discontinuity itself engineered to imprint a controlled phase shift on transmitted channels.

\subsection{P-MM-6. Coarse-graining window}

All derivations operate in the window
\[
\ell_P \ll a \ll \lambda \ll L,
\]
where $\ell_P$ is the substrate cutoff (smallest substrate length scale), $a$ is the microstructure period (engineered), $\lambda$ is the operating wavelength, and $L$ is the device scale.

\textbf{Within this window:}
\begin{itemize}[noitemsep]
\item Substrate discreteness is invisible at the wavelength scale: $\ell_P \ll \lambda$ ensures that the coarse-grained wave equation is meaningful.
\item Microstructure averaging is well-defined: $a \ll \lambda$ ensures that the two-scale expansion converges and the cell-problem corrector is well-posed.
\item Device-scale envelopes are slow: $\lambda \ll L$ ensures that boundary effects at device boundaries do not invalidate bulk derivations.
\end{itemize}

\textbf{Boundaries of the window.} The Bloch / bandgap regime lives at $a \sim \lambda$, adjacent to but outside the strict homogenization regime $a \ll \lambda$. The substrate cutoff $\ell_P$ is invisible throughout this walkthrough; nothing here depends on its precise value.

\section{Two-Scale Expansion and the Cell Problem}

This section derives the homogenization framework that produces $\varepsilon_{\rm eff}$ and $\mu_{\rm eff}$. The two-scale expansion is the mathematical reflection of P-MM-2 and P-MM-6.

\subsection{Setup}

Under P-MM-1, P-MM-2, P-MM-4, P-MM-6, consider the scalar form of the wave equation
\[
\partial_{x^i}\!\left[A^{ij}\!\left(x/a\right)\,\partial_{x^j}\psi(x)\right] + k_0^2\,n^2\!\left(x/a\right)\,\psi(x) = 0,
\]
where $A^{ij}(y)$ and $n^2(y)$ are $Y$-periodic functions of the fast coordinate $y = x/a$. The coefficient $A^{ij}(y)$ is positive-definite and bounded:
\[
c_- |\xi|^2 \le A^{ij}(y)\,\xi_i\,\xi_j \le c_+ |\xi|^2, \qquad 0 < c_- \le c_+ < \infty,
\]
for all $\xi \in \mathbb{R}^3$ and all $y \in Y$. The ratio $a/\lambda$ is small; we expand the equation in this parameter.

\subsection{Two-scale lift}

Lift $\psi(x)$ to a two-scale function $\tilde\psi(X, y)$, treating $X = x$ (slow) and $y = x/a$ (fast) as independent variables. The chain rule gives
\[
\frac{\partial}{\partial x^i} = \frac{\partial}{\partial X^i} + \frac{1}{a}\frac{\partial}{\partial y^i}.
\]
This is the \textbf{derivative splitting} rule of two-scale analysis. It expresses the fact that a spatial gradient $\partial/\partial x$ has two contributions: a slow contribution from the device-scale envelope $\partial/\partial X$ and a fast contribution from the cell-scale modulation $a^{-1}\partial/\partial y$.

Lift $\psi$ to an asymptotic series in $a$:
\[
\tilde\psi(X, y) = \psi_0(X, y) + a\psi_1(X, y) + a^2\psi_2(X, y) + a^3\psi_3(X, y) + \cdots,
\]
with each $\psi_k(X, y)$ smooth in $X$ and $Y$-periodic in $y$. Substitute into the wave equation and collect terms by power of $a$.

\subsection{Order-by-order equations}

The wave equation, in lifted form, reads
\[
\left(\partial_{X^i} + a^{-1}\partial_{y^i}\right)\!\left[A^{ij}(y)\left(\partial_{X^j} + a^{-1}\partial_{y^j}\right)\tilde\psi\right] + k_0^2 n^2(y)\tilde\psi = 0.
\]
Expand the bracket and collect powers of $a$.

\textbf{Order $a^{-2}$:}
\[
\partial_{y^i}\!\left[A^{ij}(y)\partial_{y^j}\psi_0(X, y)\right] = 0.
\]
Multiply by $\psi_0$ and integrate over $Y$; integrate by parts using $Y$-periodicity (boundary terms cancel between opposite faces of $Y$):
\[
\int_Y A^{ij}(y)\partial_{y^i}\psi_0\partial_{y^j}\psi_0\, d^3y = 0.
\]
By positive-definiteness of $A^{ij}$ (P-MM-4), the integrand is non-negative and vanishes only if $\partial_{y^j}\psi_0 = 0$. Therefore $\psi_0(X, y) = \psi_0(X)$ is independent of $y$. \textbf{The leading-order field has no microstructure dependence.}

\textbf{Order $a^{-1}$:}
\[
\partial_{y^i}\!\left[A^{ij}(y)\partial_{y^j}\psi_1(X, y)\right] = -\partial_{y^i}\!\left[A^{ij}(y)\partial_{X^j}\psi_0(X)\right].
\]
The right-hand side is linear in $\partial_{X^j}\psi_0$. By linearity of the elliptic equation in $\psi_1$, the solution has the form
\[
\psi_1(X, y) = \chi^j(y)\partial_{X^j}\psi_0(X) + \tilde\psi_1(X),
\]
where the \textbf{corrector field} $\chi^j(y)$ satisfies the \textbf{cell problem}:
\[
\boxed{\partial_{y^i}\!\left[A^{ij}(y)\partial_{y^j}\chi^k(y)\right] = -\partial_{y^i} A^{ik}(y), \qquad \chi^k \text{ is } Y\text{-periodic, with } \langle\chi^k\rangle = 0.}
\]

The cell problem is a linear elliptic boundary-value problem on $Y$ with periodic boundary conditions. By the Fredholm alternative for elliptic operators with periodic boundary conditions, it has a unique solution up to additive constants; the zero-mean condition $\langle\chi^k\rangle = 0$ fixes the constant.

\textbf{Substrate-level meaning of the corrector field.} $\chi^k(y)$ is the local accommodation pattern of the substrate response to a unit slow gradient in the $k$-direction. The substrate's rule-type structure adjusts within each unit cell to absorb the slow gradient; $\chi^k$ records that adjustment.

\subsection{The averaging operator}

Define the cell-averaging operator
\[
\langle f \rangle \equiv \frac{1}{|Y|}\int_Y f(y)\, d^3y.
\]
Two basic properties follow from $Y$-periodicity:

\begin{enumerate}[noitemsep]
\item \textbf{Vanishing of pure $y$-divergences.} For any $Y$-periodic vector field $V^i(y)$,
\[
\langle\partial_{y^i} V^i\rangle = \frac{1}{|Y|}\int_Y \partial_{y^i} V^i\, d^3y = \frac{1}{|Y|}\oint_{\partial Y} V^i n_i\, dS = 0,
\]
because the boundary integral on $\partial Y$ vanishes by periodicity (opposite faces cancel).
\item \textbf{Commutativity with $X$-derivatives.} Since $X$ and $y$ are independent in the lifted picture,
\[
\langle\partial_{X^i} f(X, y)\rangle = \partial_{X^i}\langle f(X, y)\rangle.
\]
\end{enumerate}

These two properties drive the homogenization computation.

\subsection{Order-$a^0$ equation and effective tensor}

\textbf{Order $a^0$:}
\[
\partial_{X^i}\!\left[A^{ij}\partial_{X^j}\psi_0\right] + \partial_{X^i}\!\left[A^{ij}\partial_{y^j}\psi_1\right] + \partial_{y^i}\!\left[A^{ij}\partial_{X^j}\psi_1\right] + \partial_{y^i}\!\left[A^{ij}\partial_{y^j}\psi_2\right] + k_0^2 n^2\psi_0 = 0.
\]

Apply the averaging operator. The pure $\partial_{y^i}[\cdots]$ terms (third and fourth terms) vanish by property 1. The second term, using $\psi_1 = \chi^j\partial_{X^j}\psi_0$ and property 2, becomes $\partial_{X^i}\langle A^{ij}\partial_{y^j}\chi^k\rangle\partial_{X^k}\psi_0$. The first term becomes $\partial_{X^i}\langle A^{ij}\rangle\partial_{X^j}\psi_0$. The fifth term becomes $k_0^2\langle n^2\rangle\psi_0$.

Combining:
\[
\partial_{X^i}\!\left[\big(\langle A^{ik}\rangle + \langle A^{ij}\partial_{y^j}\chi^k\rangle\big)\partial_{X^k}\psi_0\right] + k_0^2\langle n^2\rangle\psi_0 = 0.
\]

This is the \textbf{homogenized wave equation} at leading order:
\[
\partial_{X^i}\!\left[A^{*ik}\partial_{X^k}\psi_0\right] + k_0^2 n^{*2}\psi_0 = 0,
\]
with the \textbf{effective coefficient}
\[
\boxed{A^{*ik} = \langle A^{ik}\rangle + \langle A^{ij}\partial_{y^j}\chi^k\rangle, \qquad n^{*2} = \langle n^2\rangle.}
\]

The effective coefficient is the cell-averaged coefficient \emph{plus} a correction from the cell-problem solution. The correction captures the local accommodation of the field to the microstructure: as the field threads through the unit cell, it deforms to follow the substrate's rule-type response, and the deformation feeds back into the cell-averaged flux.

\subsection{Voigt-Reuss bounds}

The corrector correction is non-positive in the energy norm. Standard variational arguments on the cell problem give
\[
\langle (A^{-1})^{ij}\rangle^{-1} \le A^{*ij} \le \langle A^{ij}\rangle
\]
as quadratic forms on $\mathbb{R}^3$. The lower (Reuss) bound is the harmonic mean of $A^{ij}$ over the cell; the upper (Voigt) bound is the arithmetic mean. The effective coefficient lives between these bounds and saturates the upper bound only when $A^{ij}$ is constant on $Y$ (no microstructure, trivial homogenization).

\subsection{Worked example: layered substrate}

For a one-dimensional layered substrate with $A(y) = A_1$ on layer 1 (volume fraction $f$) and $A(y) = A_2$ on layer 2 (volume fraction $1 - f$), the cell problem in 1D reduces to
\[
\frac{d}{dy}\!\left[A(y)\left(\frac{d\chi}{dy} + 1\right)\right] = 0,
\]
so $A(y)(d\chi/dy + 1) = C$ (constant on $Y$). Solving:
\[
\frac{d\chi}{dy} + 1 = \frac{C}{A(y)}, \qquad \chi \text{ periodic} \Rightarrow \langle d\chi/dy\rangle = 0 \Rightarrow 1 = C\langle 1/A\rangle.
\]
So $C = \langle 1/A\rangle^{-1}$, and the effective coefficient parallel to the layering is
\[
A^*_\parallel = \langle A(d\chi/dy + 1)\rangle = \langle C\rangle = \langle 1/A\rangle^{-1},
\]
the harmonic mean. Perpendicular to the layering, $A^*_\perp = \langle A\rangle$, the arithmetic mean. The microstructure thus produces anisotropy even from isotropic constituents --- a substrate-level FORCED feature of any non-trivial periodic microstructure.

\section{Effective Constitutive Tensors and Negative Index}

\subsection{From scalar form to Maxwell}

For Maxwell's equations in a non-magnetic medium ($\mu = \mu_0$), the TE polarization wave equation in 2D is
\[
\partial_i\!\left[\frac{1}{\varepsilon(x)/\varepsilon_0}\partial_i H_z\right] + \frac{\omega^2}{c^2} H_z = 0,
\]
identifying $A^{ij} \leftrightarrow (\varepsilon/\varepsilon_0)^{-1}\delta^{ij}$. For TM polarization,
\[
\partial_i\!\left[\frac{1}{\mu(x)/\mu_0}\partial_i E_z\right] + \frac{\omega^2}{c^2} E_z = 0,
\]
identifying $A^{ij} \leftrightarrow (\mu/\mu_0)^{-1}\delta^{ij}$. Applying \S3 in each polarization gives effective constitutive tensors.

\subsection{Effective constitutive tensors}

In full vector form, for an anisotropic periodic substrate with $\varepsilon^{ij}(y)$ and $\mu^{ij}(y)$ both $Y$-periodic, the homogenization theorem yields
\[
\boxed{\varepsilon_{\rm eff}^{ij} = \langle\varepsilon^{ij}\rangle + \langle\varepsilon^{ik}\partial_{y^k}\chi_E^j\rangle, \qquad (\mu_{\rm eff}^{-1})^{ij} = \langle(\mu^{-1})^{ij}\rangle + \langle(\mu^{-1})^{ik}\partial_{y^k}\chi_M^j\rangle,}
\]
with $\chi_E^j(y)$ and $\chi_M^j(y)$ the electric and magnetic correctors solving their respective cell problems.

\textbf{Anisotropy from microstructure.} Even if $\varepsilon^{ij}(y) = \varepsilon(y)\delta^{ij}$ and $\mu^{ij}(y) = \mu(y)\delta^{ij}$ are locally isotropic, the corrector solutions $\chi_E^j$, $\chi_M^j$ generally produce anisotropy in $\varepsilon_{\rm eff}^{ij}$ and $\mu_{\rm eff}^{ij}$. The principal axes of the effective tensors are aligned with the symmetry axes of the microstructure.

\subsection{Substrate-level meaning of $\varepsilon$ and $\mu$}

In the ED picture, $\varepsilon^{ij}(y)$ is the substrate's local rule-type polarizability: the coarse-grained tendency for ED-flow to align under an applied electric tension. $(\mu^{-1})^{ij}(y)$ is the local rule-type inverse circulation susceptibility: the coarse-grained resistance to closed-loop ED-flow under an applied magnetic tension. Neither is a fundamental field; both are coarse-grained labels for rule-type response.

The corrector fields $\chi_E^j(y)$ and $\chi_M^j(y)$ are the local accommodation patterns of the substrate response. When a slow gradient $\partial_{X^j}\psi_0$ threads through a unit cell, the substrate's rule-type structure adjusts within the cell; $\chi^j(y)$ records the spatial pattern of that adjustment, and the cell-averaged adjustment contributes to the effective tensor.

\subsection{Wire-array microstructure $\to$ Drude $\varepsilon_{\rm eff}(\omega)$}

Consider a square array of thin conducting wires of radius $r_0$ and lattice constant $a$, with $r_0 \ll a \ll \lambda$. Standard analysis (Pendry 1996) of the cell problem with electric field along the wire axis identifies the wires as supporting a longitudinal collective oscillation with effective plasma frequency
\[
\omega_p^2 = \frac{2\pi c^2}{a^2 \ln(a/r_0)}.
\]
The wire-array effective dielectric function takes the Drude form
\[
\boxed{\varepsilon_{\rm eff}(\omega) = \varepsilon_\infty\left(1 - \frac{\omega_p^2}{\omega^2}\right).}
\]
Below $\omega_p$, $\varepsilon_{\rm eff}(\omega) < 0$. Above $\omega_p$, $\varepsilon_{\rm eff}(\omega) > 0$ and approaches $\varepsilon_\infty$ at high frequency.

\textbf{Substrate-level meaning.} The wire array carries a collective rule-type oscillation: ED-flow along the wires is constrained by the wire geometry to a single longitudinal mode, with effective plasma frequency set by the geometric line density and the cell logarithm. Below this frequency, the substrate's longitudinal response is \emph{opposite} to the applied tension --- the engineered rule-type collective mode dominates over any local polarizability --- producing negative effective $\varepsilon$.

\subsection{Split-ring microstructure $\to$ Lorentz $\mu_{\rm eff}(\omega)$}

Consider an array of split-ring resonators (SRRs) --- pairs of concentric metallic loops with a gap, embedded in a host dielectric --- with lattice constant $a$, resonance frequency $\omega_0$, oscillator strength $F$ (related to the ring's filling fraction), and damping rate $\gamma$. Standard analysis of the cell problem for an applied magnetic field perpendicular to the rings (Pendry 1999) gives the Lorentz form
\[
\boxed{\mu_{\rm eff}(\omega) = 1 - \frac{F\omega^2}{\omega^2 - \omega_0^2 + i\gamma\omega}.}
\]
For $\omega$ in a narrow band immediately above $\omega_0$, $\mathrm{Re}\,\mu_{\rm eff}(\omega) < 0$.

\textbf{Substrate-level meaning.} Each SRR is a closed-loop rule-type resonator: ED-flow circulating around the ring constitutes a coarse-grained magnetic moment, whose natural frequency $\omega_0$ is set by the ring's geometric inductance and capacitance (gap capacitance). Driven above $\omega_0$, the rule-type circulation lags by more than $\pi/2$ and opposes the applied magnetic tension, producing negative effective $\mu$.

\subsection{The doubly-negative band}

Engineering both wire arrays and SRRs into the same substrate produces a frequency band where both $\varepsilon_r(\omega) < 0$ and $\mu_r(\omega) < 0$. In this band, $\varepsilon_r\mu_r > 0$, so the wave equation
\[
\nabla^2\psi + k_0^2\varepsilon_r\mu_r\psi = 0
\]
admits propagating solutions $\psi = e^{i\mathbf{k}\cdot\mathbf{x}}$ with $|\mathbf{k}|^2 = k_0^2\varepsilon_r\mu_r > 0$.

\subsection{Negative-index branch FORCED}

The refractive index satisfies $n_{\rm eff}^2 = \varepsilon_r\mu_r$, so $n_{\rm eff} = \pm\sqrt{\varepsilon_r\mu_r}$. The branch is determined by causal outgoing-wave selection.

\textbf{Argument.} In the doubly-negative band, the relation between $\mathbf{D}$ and $\mathbf{E}$ is $\mathbf{D} = \varepsilon_r\varepsilon_0\mathbf{E}$ with $\varepsilon_r < 0$, reversing direction. Similarly, $\mathbf{B} = \mu_r\mu_0\mathbf{H}$ with $\mu_r < 0$ reverses direction. The Poynting vector
\[
\mathbf{S} = \mathbf{E}\times\mathbf{H}
\]
remains in the direction of energy outflow (causality, P-MM-4). For a plane wave, the phase velocity $\mathbf{v}_p = (\omega/|\mathbf{k}|^2)\mathbf{k}$ is parallel to $\mathbf{k}$. Combining Maxwell's curl equations $\nabla\times\mathbf{E} = -\partial_t\mathbf{B}$ and $\nabla\times\mathbf{H} = \partial_t\mathbf{D}$:
\[
\mathbf{k}\times\mathbf{E} = \omega\mu_r\mu_0\mathbf{H}, \qquad \mathbf{k}\times\mathbf{H} = -\omega\varepsilon_r\varepsilon_0\mathbf{E}.
\]
Both signs flip when $\varepsilon_r, \mu_r < 0$. The triad $(\mathbf{E}, \mathbf{H}, \mathbf{k})$ becomes left-handed instead of right-handed. The Poynting vector $\mathbf{S} = \mathbf{E}\times\mathbf{H}$ now points \emph{opposite} to $\mathbf{k}$.

For an outgoing wave with $\mathbf{S}$ pointing away from the source, $\mathbf{k}$ points toward the source. The refractive index $n_{\rm eff} = c|\mathbf{k}|/\omega$ is conventionally signed by the direction of $\mathbf{k}$ relative to the outward normal. With $\mathbf{k}$ inward and $\mathbf{S}$ outward, this is precisely the negative-index branch:
\[
\boxed{n_{\rm eff}(\omega) = -\sqrt{\varepsilon_r(\omega)\mu_r(\omega)} \qquad \text{when } \varepsilon_r(\omega), \mu_r(\omega) < 0.}
\]
This is the Pendry 2000 negative refractive index. The choice is FORCED by causal outgoing-wave selection given the doubly-negative constitutive response; no separate postulate is required.

\subsection{Substrate-level meaning of negative index}

The substrate's rule-type structure in the doubly-negative band is engineered such that channels propagate with phase advancing in the direction \emph{opposite} to energy flow. In substrate terms: the engineered combination of collective rule-type oscillation (wire-array Drude response) and resonant rule-type circulation (SRR Lorentz response) produces a coarse-grained substrate whose channels carry phase upstream relative to their energy flow. Negative refraction at an interface --- bending to the same side of the normal as the incident wave --- is the direct geometric consequence at the interface boundary.

\subsection{Bandwidth and dispersion bounds}

The wire-array and SRR responses are both intrinsically dispersive: by Kramers-Kronig relations on causal linear response, $\mathrm{Im}\,\varepsilon(\omega) \neq 0$ over any band of finite width implies $\mathrm{Re}\,\varepsilon(\omega)$ varies with $\omega$. The doubly-negative band, where both $\mathrm{Re}\,\varepsilon$ and $\mathrm{Re}\,\mu$ are negative, is intrinsically narrow ($\Delta\omega/\omega \sim 5$--$20\%$ in typical microwave metamaterials). Lossless negative-index media over arbitrarily broad bandwidth are forbidden by causality.

\section{Substrate Deformation, Effective Metric, and the Pendry Cloak}

\subsection{Setup}

Under P-MM-1, P-MM-3, P-MM-4, P-MM-6, consider a slow substrate deformation $x'^i = f^i(x)$ with Jacobian
\[
J^i{}_j = \frac{\partial x'^i}{\partial x^j}, \qquad (J^{-1})^k{}_l = \frac{\partial x^k}{\partial x'^l}, \qquad \det J > 0.
\]
The undeformed substrate carries an isotropic vacuum rule-type response $\varepsilon^{ij}(x) = \varepsilon_0\delta^{ij}$, $\mu^{ij}(x) = \mu_0\delta^{ij}$. The deformation maps this to a coarse-grained anisotropic medium with transformed constitutive tensors derived below.

\subsection{Rule-type deformation tensor}

The substrate's local rule-type orientation is a rank-2 contravariant tensor structure $D^{ij}(x)$ at each point. Under deformation, $D^{ij}$ transports covariantly as
\[
D'^{ij}(x') = \frac{1}{\det J} J^i{}_k J^j{}_l D^{kl}(x).
\]
The factor $1/\det J$ accounts for the volume change under deformation (density of weight one). The Jacobian factors transport the tensor indices into the deformed frame.

\textbf{Conjugate tensor.} The conjugate rank-2 covariant tensor structure $D_{ij}(x)$ transforms as
\[
D'_{ij}(x') = (\det J)(J^{-1})^k{}_i (J^{-1})^l{}_j D_{kl}(x).
\]

\subsection{Constitutive transformation}

The substrate's coarse-grained electric and magnetic responses are rank-2 contravariant tensors. Applying the deformation rule of \S5.2:
\[
\boxed{\varepsilon'^{ij}(x') = \frac{1}{\det J} J^i{}_k J^j{}_l \varepsilon^{kl}(x), \qquad \mu'^{ij}(x') = \frac{1}{\det J} J^i{}_k J^j{}_l \mu^{kl}(x).}
\]

This is the \textbf{transformation-optics constitutive law}. It is a substrate-level FORCED statement: given the deformation $f$, the transformed constitutive tensors are computed by applying the Jacobian densities. No separate posit is required.

\textbf{Special case: vacuum starting medium.} If the undeformed substrate is vacuum ($\varepsilon^{ij} = \varepsilon_0\delta^{ij}$, $\mu^{ij} = \mu_0\delta^{ij}$), then
\[
\varepsilon'^{ij}(x') = \frac{\varepsilon_0}{\det J}(JJ^T)^{ij}, \qquad \mu'^{ij}(x') = \frac{\mu_0}{\det J}(JJ^T)^{ij}.
\]
The impedance $\sqrt{\mu'/\varepsilon'}$ equals $\sqrt{\mu_0/\varepsilon_0} = Z_0$ everywhere --- the transformed medium is \textbf{impedance-matched to vacuum at every point}, eliminating reflections from interior surfaces.

\subsection{Effective metric}

The wave equation in the deformed substrate, with the transformed constitutive tensors, can be rewritten in covariant Laplace-Beltrami form. Consider Maxwell's equations in a transformed medium; using the transformation rule of \S5.3 in the scalar-form wave equation gives
\[
\partial_{x'^i}\!\left[\frac{(JJ^T)^{ij}}{\det J}\partial_{x'^j}\psi\right] + k_0^2\psi = 0.
\]

Identify
\[
\boxed{g^{ij}(x') = \frac{(JJ^T)^{ij}}{\det J}, \qquad g_{ij}(x') = (J^{-1})^k{}_i (J^{-1})^l{}_j \delta_{kl}, \qquad \sqrt{|g|} = \frac{1}{\det J}.}
\]

The wave equation becomes
\[
\frac{1}{\sqrt{|g|}}\partial_{x'^i}\!\left[\sqrt{|g|}\,g^{ij}\partial_{x'^j}\psi\right] + k_0^2\psi = 0.
\]

This is the \textbf{covariant Laplace-Beltrami wave equation} on a Riemannian manifold with metric $g_{ij}$.

\textbf{Crucial scope note.} The metric $g_{ij}$ is a coarse-grained rule-type metric on the substrate: it encodes the coarse-grained rule-type geometry seen by the propagating field. It is \emph{not} spacetime curvature. The substrate-level operation that produced it is an engineered displacement (P-MM-3), not a gravitational effect. The identification $g_{ij} \leftrightarrow$ rule-type metric is structural, not analogical: the same mathematical object appears in both contexts because both contexts describe wave propagation through a curved geometry, but the underlying physics is different.

\subsection{Eikonal limit: rays as geodesics}

In the eikonal limit $\lambda \to 0$, the wave equation reduces to the eikonal equation
\[
g^{ij}\partial_i S\,\partial_j S = n^2,
\]
where $S(x')$ is the phase function and $n$ is the local refractive index. The solutions are geodesics of the effective metric $g_{ij}$. \textbf{Light in a transformation-optics medium follows geodesics of the engineered substrate metric.} This is the operational statement that drives cloaking, lensing, and waveguiding in the transformation-optics regime.

\subsection{Pendry spherical cloak: deformation}

Choose the radial map
\[
r' = R_1 + \alpha r, \qquad \alpha = \frac{R_2 - R_1}{R_2}, \qquad 0 \le r \le R_2,
\]
in spherical coordinates $(r, \theta, \phi)$, with angular coordinates unchanged. The map sends the open ball $r < R_2$ onto the shell $R_1 < r' < R_2$. The interior point $r = 0$ maps to the inner shell $r' = R_1$; the outer boundary $r = R_2$ maps to itself.

The Jacobian in spherical coordinates has principal components:
\begin{itemize}[noitemsep]
\item Radial: $\partial r'/\partial r = \alpha$.
\item Polar angular: $r'/r$.
\item Azimuthal angular: $(r'/r)\sin\theta/\sin\theta = r'/r$.
\end{itemize}

Note $r = (r' - R_1)/\alpha$. So $r'/r = \alpha r'/(r' - R_1)$.

\subsection{Pendry cloak: constitutive tensors}

Applying the constitutive transformation of \S5.3 in spherical coordinates with $\det J = \alpha(r'/r)^2 = \alpha(\alpha r'/(r' - R_1))^2$:
\[
\det J = \alpha^3 \frac{r'^2}{(r' - R_1)^2}.
\]

The radial constitutive component:
\[
\frac{\varepsilon_{\hat r'}}{\varepsilon_0} = \frac{J_{rr}^2}{\det J} \cdot \frac{1}{(g_{\rm sph})_{rr}} \;\Rightarrow\; \frac{\varepsilon_{\hat r'}}{\varepsilon_0} = \beta\frac{(r' - R_1)^2}{r'^2},
\]
where $\beta = R_2/(R_2 - R_1) = 1/\alpha$ (with appropriate handling of the spherical metric factors).

The angular constitutive components:
\[
\frac{\varepsilon_{\hat\theta'}}{\varepsilon_0} = \frac{\varepsilon_{\hat\phi'}}{\varepsilon_0} = \beta.
\]

The same expressions hold for $\mu/\mu_0$ --- the cloak is impedance-matched to vacuum at every point. In boxed form:
\[
\boxed{\frac{\varepsilon_{\hat r'}}{\varepsilon_0} = \beta\frac{(r' - R_1)^2}{r'^2}, \qquad \frac{\varepsilon_{\hat\theta'}}{\varepsilon_0} = \frac{\varepsilon_{\hat\phi'}}{\varepsilon_0} = \beta, \qquad \beta = \frac{R_2}{R_2 - R_1}.}
\]

These are the \textbf{Pendry spherical cloak constitutive tensors} (Pendry 2006).

\subsection{Effective metric of the cloak}

The effective metric components, applying \S5.4 to the cloak deformation:
\[
g_{\hat r'\hat r'} = \beta^2\frac{(r' - R_1)^2}{r'^2}, \qquad g_{\hat\theta'\hat\theta'} = g_{\hat\phi'\hat\phi'} = \beta^2.
\]

At $r' = R_1$, the radial metric component vanishes: $g_{\hat r'\hat r'}(r' = R_1) = 0$.

\subsection{Substrate-level meaning of invisibility}

The vanishing of $g_{\hat r'\hat r'}$ at $r' = R_1$ has the substrate-level reading that \textbf{rule-type channels carrying field amplitude cannot enter the cloaked interior}. The radial distance, in the rule-type metric, becomes singular at the inner shell: any geodesic of $g_{ij}$ that attempts to cross $r' = R_1$ would require infinite proper length. The interior is \emph{topologically excluded} from the substrate seen by the propagating field.

This is the structural meaning of invisibility:
\begin{itemize}[noitemsep]
\item Not absorption (no energy loss in the cloak shell).
\item Not reflection (impedance matching eliminates reflections).
\item Not concealment (no hiding behind material screens).
\item \textbf{Topological exclusion}: the cloaked region is removed from the rule-type geometry seen by the propagating channel.
\end{itemize}

Energy flow follows geodesics of $g_{ij}$, which detour smoothly around the cloaked interior and exit the outer shell $r' = R_2$ as if traveling through vacuum. An observer outside the cloak sees the field profile of vacuum propagation; the cloaked object is invisible.

\subsection{General transformation optics}

The Pendry cloak is one instance of the general transformation-optics construction. Any smooth deformation $f: x \mapsto x'$ with $J$ and $J^{-1}$ bounded and $\det J > 0$ produces a substrate-level effective medium with constitutive tensors given by \S5.3 and effective metric given by \S5.4. The geodesics of the effective metric are the engineered light rays.

\textbf{Physical realizability conditions:}
\begin{itemize}[noitemsep]
\item Smoothness of $f$: $J, J^{-1}$ smooth and bounded.
\item Orientation preservation: $\det J > 0$.
\item Homogenization regime: the resulting $\varepsilon', \mu'$ must be achievable by some periodic microstructure with $a \ll \lambda$ at the operating frequency.
\item Absence of unintentional singularities: $\varepsilon'$ and $\mu'$ singularities only at designed locations.
\item Bandwidth: dispersion of the implementing microstructure restricts the design to a finite frequency band.
\end{itemize}

The Pendry cloak saturates condition 4 deliberately at $r' = R_1$; this singularity is the price of perfect cloaking, and it is what makes the perfect cloak intrinsically bandwidth-limited and physically realizable only as an approximation.

\section{Metasurfaces and the Generalized Snell's Law}

\subsection{Codimension-1 discontinuity}

Under P-MM-5, P-MM-6, consider a surface $\Sigma$ at $z = 0$ across which the substrate's rule-type structure changes discontinuously. The bulk regions $z < 0$ and $z > 0$ are described by the coarse-grained wave equation of P-MM-4 with refractive indices $n_1$ and $n_2$ respectively. The discontinuity at $\Sigma$ imposes jump conditions on the wave field.

\subsection{Passive interface: standard Snell's law}

For a passive interface (no engineered surface phase profile, no surface polarization, no surface currents), the standard tangential-field continuity holds:
\[
\mathbf{E}_\parallel(r_\parallel, 0^+) = \mathbf{E}_\parallel(r_\parallel, 0^-), \qquad \mathbf{H}_\parallel(r_\parallel, 0^+) = \mathbf{H}_\parallel(r_\parallel, 0^-).
\]
For a scalar wave $\psi$ (component of $\mathbf{E}$ or $\mathbf{H}$ parallel to $\Sigma$):
\[
\psi(r_\parallel, 0^+) = \psi(r_\parallel, 0^-).
\]

Take a plane wave on side 1: $\psi_1 = e^{i k_{\parallel,1}\cdot r_\parallel + i k_{z,1} z}$. Continuity at $z = 0$ forces the tangential wavevector to match: $k_{\parallel,2} = k_{\parallel,1}$. With $k_{\parallel,1} = n_1 k_0 \sin\theta_i$ and $k_{\parallel,2} = n_2 k_0 \sin\theta_t$, the matching condition gives
\[
n_1 \sin\theta_i = n_2 \sin\theta_t,
\]
the standard Snell's law.

\subsection{Engineered phase profile}

A metasurface is a codimension-1 surface engineered to imprint a tangential phase profile $\Phi(r_\parallel)$ on the transmitted field. The phase profile is produced by subwavelength resonant scatterers (``meta-atoms'') whose engineered geometry imprints a position-dependent phase shift. Each meta-atom is a coarse-grained rule-type discontinuity of negligible thickness; the metasurface as a whole is the spatial pattern of these meta-atoms.

The boundary condition at $\Sigma$ becomes, for a scalar wave:
\[
\boxed{\psi_2(r_\parallel, 0^+) = \psi_1(r_\parallel, 0^-)\,e^{i\Phi(r_\parallel)}.}
\]

\textbf{Substrate-level meaning.} The metasurface is a codimension-1 engineered rule-type discontinuity. The tangential phase profile $\Phi(r_\parallel)$ is the integrated rule-type phase imprinted on a channel as it crosses $\Sigma$. The function $\Phi$ is engineered by the meta-atom geometry: each meta-atom resonance produces a controlled phase shift, and the spatial arrangement of meta-atoms produces the spatial profile.

\subsection{Tangential momentum kick}

Take an incident plane wave on side 1: $\psi_1 = e^{i k_{\parallel,1}\cdot r_\parallel + i k_{z,1} z}$. At $z = 0^-$, the field is $\psi_1(r_\parallel, 0^-) = e^{i k_{\parallel,1}\cdot r_\parallel}$. Applying the metasurface boundary condition:
\[
\psi_2(r_\parallel, 0^+) = e^{i k_{\parallel,1}\cdot r_\parallel + i\Phi(r_\parallel)}.
\]

Take the tangential gradient of the phase. For a \emph{linear} phase profile $\Phi(x) = (d\Phi/dx)x$:
\[
\psi_2(r_\parallel, 0^+) = e^{i[k_{\parallel,1} + d\Phi/dx]x},
\]
showing that the tangential wavevector on side 2 is
\[
\boxed{k_{\parallel,2} = k_{\parallel,1} + \nabla_\parallel\Phi(r_\parallel).}
\]

The metasurface imparts a \textbf{tangential momentum kick} equal to the gradient of the engineered phase profile. This is the operational signature of the metasurface: the discontinuity adds a tangential wavevector contribution that no passive interface can produce.

\subsection{Generalized Snell's law}

Substitute $k_{\parallel,1} = n_1 k_0 \sin\theta_i$, $k_{\parallel,2} = n_2 k_0 \sin\theta_t$, $k_0 = 2\pi/\lambda$ into the matching condition $k_{\parallel,2} = k_{\parallel,1} + d\Phi/dx$:
\[
n_2 k_0 \sin\theta_t = n_1 k_0 \sin\theta_i + \frac{d\Phi}{dx},
\]
or equivalently
\[
\boxed{n_1 \sin\theta_i = n_2 \sin\theta_t + \frac{\lambda}{2\pi}\frac{d\Phi}{dx}.}
\]

This is the \textbf{Capasso 2011 generalized Snell's law}. The standard Snell's law is the special case $d\Phi/dx = 0$.

\subsection{Consequences}

The generalized Snell's law permits:

\begin{itemize}[noitemsep]
\item \textbf{Anomalous refraction.} For $d\Phi/dx \neq 0$, the refraction angle $\theta_t$ is no longer determined solely by the bulk indices and incidence angle. Refraction angles unattainable from $n_1, n_2$ alone become accessible.
\item \textbf{Negative refraction at a single interface.} A sufficiently steep phase gradient $d\Phi/dx$ can produce $\theta_t$ on the same side of the normal as $\theta_i$ even with both bulk indices positive.
\item \textbf{Reflectionless beam steering.} By engineering $\Phi$, the metasurface redirects a transmitted beam to any angle within the limits of the phase gradient.
\item \textbf{Phase-only optical components.} Lenses, waveplates, holograms, and arbitrary optical functionality implemented as flat metasurfaces of subwavelength thickness.
\end{itemize}

\subsection{Limits}

The phase profile $\Phi(r_\parallel)$ must vary slowly on the scale of $\lambda$ for the boundary-condition treatment to be valid. Phase profiles with features finer than $\lambda$ are not captured by the metasurface picture alone; they enter full diffraction theory. The phase imprint is intrinsically frequency-dependent (the meta-atoms supporting $\Phi$ are resonant), so achromatic metasurfaces require multi-resonance engineering and remain bandwidth-limited.

\section{Photonic Bandgaps via Bloch Eigenmodes}

\subsection{When homogenization breaks down}

The two-scale expansion of \S3 requires $a/\lambda \to 0$. At $a \sim \lambda$, the expansion no longer converges term-by-term: the corrector field develops slow-coordinate dependence, and additional corrections of order $(a/\lambda)^2$ enter the effective coefficient. The relevant description in this regime is Bloch / Floquet analysis: instead of effective constants, one obtains a band structure $\omega_n(k)$ on the Brillouin zone.

This is the regime of \textbf{photonic crystals}, adjacent to the strict homogenization regime.

\subsection{Bloch's theorem on the periodic substrate}

Under P-MM-1, P-MM-2, P-MM-4, P-MM-6 with $a \sim \lambda$, consider the wave equation
\[
\partial_i\!\left[A^{ij}(x)\partial_j\psi\right] + \omega^2 n^2(x)\psi = 0,
\]
with $A^{ij}$ and $n^2$ periodic with lattice $\Lambda = a\mathbb{Z}^3$. Define the translation operator $T_R$ acting on functions: $T_R\psi(x) = \psi(x + R)$ for $R \in \Lambda$.

The wave operator commutes with $T_R$ for any $R \in \Lambda$. The simultaneous eigenfunctions of the wave operator and $\{T_R\}_{R\in\Lambda}$ form a complete basis. The translation operators are unitary on a suitable Hilbert space, so their eigenvalues are phases $e^{ik\cdot R}$ with $k$ a real vector. Restricting $k$ to the first Brillouin zone (BZ) --- the set of $k$ such that $|k| \le |k - G|$ for all reciprocal lattice vectors $G \in \Lambda^*$ --- uniquely labels the simultaneous eigenfunctions.

\textbf{Bloch form.} Simultaneous eigenfunctions take the form
\[
\boxed{\psi_{n, k}(x) = u_{n, k}(x)\,e^{i k\cdot x}, \qquad u_{n, k}(x + R) = u_{n, k}(x) \text{ for all } R \in \Lambda,}
\]
indexed by band index $n \in \mathbb{N}$ and crystal momentum $k$ in the first BZ.

\subsection{Band structure}

Substituting the Bloch form into the wave equation gives, after expanding the derivatives,
\[
(\partial_i + i k_i)\!\left[A^{ij}(x)(\partial_j + i k_j) u_{n, k}\right] + \omega_n^2(k) n^2(x) u_{n, k} = 0,
\]
a Hermitian generalized eigenvalue problem for $u_{n, k}$ on a single unit cell with periodic boundary conditions. At each $k$, the eigenvalues $\omega_n^2(k)$ form a discrete real sequence
\[
0 \le \omega_1^2(k) \le \omega_2^2(k) \le \omega_3^2(k) \le \cdots,
\]
indexed by band number $n$. As $k$ varies continuously over the BZ, each eigenvalue $\omega_n^2(k)$ traces out the \textbf{$n$-th band} of the band structure.

\subsection{Bandgap formation}

A \textbf{bandgap} is a frequency interval $[\omega_a, \omega_b]$ such that no value of $k$ produces an eigenvalue $\omega_n^2(k)$ in $[\omega_a^2, \omega_b^2]$ for any band $n$. At those frequencies, no propagating Bloch eigenmode exists --- the substrate forbids field propagation in the bulk.

\textbf{Mechanism.} A bandgap appears when the bands $\omega_n(k)$ and $\omega_{n+1}(k)$, in some region of the BZ, separate without crossing. The standard mechanism is constructive vs.\ destructive interference of the field across the unit cell: at certain $k$ near the BZ boundary, the spatial pattern of the lower band has nodes on the high-index regions while the higher band has nodes on the low-index regions, with the energy difference producing a gap.

For periodic dielectric substrates with sufficient index contrast (typically $\Delta\varepsilon/\varepsilon \gtrsim 0.5$ for full 3D gaps), complete bandgaps appear: frequency intervals in which no propagating mode exists in any direction. This is the \textbf{Yablonovitch photonic bandgap}.

\subsection{Substrate-level meaning}

The bandgap is a frequency window in which the periodic rule-type substrate offers no allowed coarse-grained channel. The substrate's rule-type structure, periodicized at scale $a \sim \lambda$, supports certain frequencies as propagating modes (the bands) and forbids others (the gaps). The forbidden frequencies are those for which no Bloch eigenmode can be assembled from the local rule-type response.

This is the substrate-level analog of an electronic bandgap in a crystalline solid: same Bloch machinery, applied to the substrate's coarse-grained rule-type response rather than electron wave functions.

\subsection{Relation to homogenization}

The homogenization regime $a \ll \lambda$ corresponds to the long-wavelength limit $k \to 0$ of the lowest band:
\[
\omega_1(k) \approx c_{\rm eff}|k|, \qquad k \to 0,
\]
with $c_{\rm eff}$ determined by the effective constitutive tensors of \S\S3-4. The lowest band, near $k = 0$, looks like vacuum propagation in a homogeneous effective medium.

Bandgaps live further into the BZ, at $|k| \sim \pi/a$, where homogenization is no longer valid and the full Bloch structure is needed. At these wavevectors, the wavelength $\lambda = 2\pi/k \sim 2a$ is comparable to the lattice period, and the field oscillates significantly within one cell.

The two regimes --- homogenization and Bloch / bandgap --- are not independent. They are two limits of the same substrate operation: engineered periodic microstructure under P-MM-1, P-MM-2. Homogenization gives the long-wavelength effective medium; Bloch analysis gives the full band structure including bandgaps. A photonic crystal is the same substrate operation as a homogenized metamaterial; the difference is the operating wavelength relative to the lattice period.

\subsection{Defects and waveguiding}

Local defects in the periodic substrate (a single perturbed unit cell, a missing inclusion, an embedded line of perturbed cells) support localized modes inside the bandgap. These defect states are the substrate-level mechanism for photonic-crystal waveguides, cavities, and add-drop filters. The defect modifies the local rule-type response, producing a discrete eigenmode whose frequency lies in the bulk's forbidden gap.

\section{The Unifying Substrate-Level Statement}

\subsection{Three control axes}

The three closures derived in this walkthrough --- homogenization, transformation optics, metasurface boundary conditions --- correspond to three independent engineered modulations of the substrate's coarse-grained rule-type structure:

\begin{center}
\footnotesize
\begin{tabular}{p{0.21\textwidth}p{0.16\textwidth}p{0.27\textwidth}p{0.26\textwidth}}
\toprule
\textbf{Operation} & \textbf{Substrate primitive} & \textbf{Coarse-grained result} & \textbf{Nobel-frontier instance} \\
\midrule
Smooth periodic microstructure & P-MM-1, P-MM-2 & $\varepsilon_{\rm eff}, \mu_{\rm eff}$ (and Bloch bands at $a\sim\lambda$) & Pendry-negative-index, Yablonovitch-bandgap \\
Smooth deformation & P-MM-3 & effective metric $g_{ij}$, transformed $\varepsilon, \mu$ & Pendry-cloak \\
Codim-1 discontinuity & P-MM-5 & tangential momentum kick $\nabla_\parallel\Phi$ & Capasso-metasurface \\
\bottomrule
\end{tabular}
\end{center}

\subsection{Exhaustiveness within the precursor machinery}

Within the coarse-graining window P-MM-6 and under the second-order wave equation P-MM-4, engineered modifications of the substrate enter one of three locations:

\begin{enumerate}[noitemsep]
\item The coefficient $A^{ij}(x)$ varies smoothly on a scale $a$ --- Closure 1.
\item The coordinate $x$ is the image of a deformed substrate $x = f(X')$ with $f$ smooth --- Closure 2.
\item The coefficient $A^{ij}(x)$ has a codimension-1 jump on a surface --- Closure 3.
\end{enumerate}

Higher-codimension features (lines, points) do not survive coarse-graining: they appear as point scatterers at coarse-grained scales and are not metamaterial control mechanisms. Volumetric non-smooth features cannot be averaged within the homogenization window and reduce to one of the three cases at the operating wavelength. \textbf{The three closures therefore exhaust the substrate-level wave-control space within the precursor machinery.}

\subsection{The unifying identity}

\[
\boxed{\text{Metamaterials} = \text{engineered rule-type microstructure} + \text{engineered deformation} + \text{engineered discontinuity}.}
\]

Each Nobel-frontier result is the instantiation of one axis in isolation:
\begin{itemize}[noitemsep]
\item \textbf{Yablonovitch bandgap} = periodic microstructure at $a \sim \lambda$.
\item \textbf{Pendry negative index} = doubly-resonant periodic microstructure at $a \ll \lambda$.
\item \textbf{Pendry cloak} = smooth deformation with FORCED constitutive transformation.
\item \textbf{Capasso metasurface} = engineered codim-1 phase discontinuity.
\end{itemize}

Practical metamaterials generally compose all three axes: a transformation-optics device implemented by a graded homogenized medium with a metasurface termination, for example.

\subsection{The unified substrate-level wave equation}

The most general coarse-grained wave equation in the precursor machinery is
\[
\frac{1}{\sqrt{|g|}}\partial_i\!\left[\sqrt{|g|}\,g^{ij}(\varepsilon_{\rm eff}^{-1})^{jk}\partial_k\psi\right] + k_0^2\mu_{\rm eff}\psi = 0,
\]
subject to jump conditions on engineered codim-1 surfaces $\Sigma$:
\[
\psi(r_\parallel, 0^+) = \psi(r_\parallel, 0^-)\,e^{i\Phi(r_\parallel)}.
\]

This is the \textbf{unified substrate-level wave equation of the Arc}. The three closures are its three independent control knobs: $\varepsilon_{\rm eff}, \mu_{\rm eff}$ (Closure 1), $g_{ij}$ (Closure 2), and $\Phi$ on $\Sigma$ (Closure 3). All three knobs derive from the same substrate operations (P-MM-1, P-MM-3, P-MM-5) coarse-grained under P-MM-6 and propagated under P-MM-4.

\subsection{Scope and limits}

The closures hold under the following bounds:

\begin{itemize}[noitemsep]
\item \textbf{Homogenization regime:} $a/\lambda \ll 1$. Beyond $a \sim \lambda/3$ the leading expansion is qualitative only; Bloch analysis takes over.
\item \textbf{Negative-index bandwidth:} $\varepsilon_r < 0$ and $\mu_r < 0$ overlap over a band intrinsically narrowed by Kramers-Kronig. Lossless negative-index media over arbitrary bandwidth are forbidden.
\item \textbf{Transformation optics smoothness:} $J$ and $J^{-1}$ bounded, $\det J > 0$. Singular maps (Pendry cloak at $r' = R_1$) are intrinsically dispersive and bandwidth-limited.
\item \textbf{Metasurface phase variation:} $\Phi(r_\parallel)$ slow on $\lambda$. Sub-$\lambda$ phase features require full diffraction theory.
\item \textbf{Universal:} energy conservation, causality (Kramers-Kronig), Sommerfeld-Brillouin signal-front bound $v_{\rm front} \le c$, bandwidth-response trade-off across all axes.
\end{itemize}

The effective metric $g_{ij}$ of \S5 is a coarse-grained rule-type metric on the substrate --- not spacetime curvature. The substrate-level operation that produces it is engineered displacement, not gravitational physics. The identification is structural, not analogical.

\section{FORCED / INHERITED / OPEN Accounting}

\subsection{FORCED at substrate level}

The following results are FORCED by the substrate primitives P-MM-1 through P-MM-6 and the derivations of this walkthrough. No additional posit is required.

\begin{itemize}[noitemsep]
\item Two-scale lift and derivative splitting (\S3.2) --- FORCED from P-MM-2, P-MM-6.
\item Order-by-order equations (\S3.3) --- FORCED.
\item Leading-order $\psi_0(X, y) = \psi_0(X)$ independence from fast variable (\S3.3) --- FORCED by positive-definiteness of $A^{ij}$.
\item Cell problem with $Y$-periodicity and zero-mean corrector (\S3.3) --- FORCED.
\item Cell-averaging operator properties (vanishing of pure $y$-divergences, commutativity with $X$-derivatives) (\S3.4) --- FORCED.
\item Effective-coefficient formula $A^{*ik} = \langle A^{ik}\rangle + \langle A^{ij}\partial_{y^j}\chi^k\rangle$ (\S3.5) --- FORCED.
\item Voigt-Reuss bounds (\S3.6) --- FORCED.
\item Anisotropy from microstructure (\S3.7, \S4.2) --- FORCED FORM.
\item Form of effective constitutive tensors $\varepsilon_{\rm eff}^{ij}$, $\mu_{\rm eff}^{ij}$ (\S4.2) --- FORCED FORM.
\item Negative-index branch choice $n = -\sqrt{\varepsilon_r\mu_r}$ in the doubly-negative band (\S4.7) --- FORCED by causal outgoing-wave selection.
\item Constitutive transformation under deformation $\varepsilon'^{ij} = (\det J)^{-1} J^i{}_k J^j{}_l \varepsilon^{kl}$ (\S5.3) --- FORCED from P-MM-3 + P-MM-4 + P-MM-6.
\item Impedance matching to vacuum under deformation from vacuum starting medium (\S5.3) --- FORCED.
\item Effective metric $g_{ij} = (J^{-1})^k{}_i (J^{-1})^l{}_j \delta_{kl}$ (\S5.4) --- FORCED.
\item Light rays as geodesics of $g_{ij}$ in eikonal limit (\S5.5) --- FORCED.
\item Pendry cloak constitutive tensors (\S5.7) --- FORCED FORM from the deformation $r' = R_1 + \alpha r$.
\item Topological exclusion of cloaked interior (\S5.9) --- FORCED by vanishing of $g_{\hat r'\hat r'}$ at $r' = R_1$.
\item Metasurface phase-imprinting boundary condition $\psi_2 = \psi_1 e^{i\Phi}$ (\S6.3) --- FORCED from P-MM-5 + P-MM-6.
\item Tangential momentum kick $k_{\parallel,2} = k_{\parallel,1} + \nabla_\parallel\Phi$ (\S6.4) --- FORCED.
\item Generalized Snell's law (\S6.5) --- FORCED.
\item Bloch form for periodic substrates (\S7.2) --- FORCED from P-MM-1, P-MM-2, P-MM-4.
\item Band structure as Hermitian eigenvalue problem on the unit cell (\S7.3) --- FORCED.
\item Photonic-bandgap mechanism as eigenvalue gaps in the Bloch problem (\S7.4) --- FORCED FORM.
\item Continuity of homogenization regime with Bloch regime through long-wavelength limit of lowest band (\S7.6) --- FORCED.
\item Three-axis exhaustiveness within P-MM-4 + P-MM-6 (\S8.2) --- FORCED.
\item Unifying identity (\S8.3) --- FORCED.
\item Unified substrate-level wave equation (\S8.4) --- FORCED FORM.
\end{itemize}

\subsection{INHERITED (form FORCED, values from microstructure or geometry choice)}

The following take their \emph{form} from the FORCED results above, but their numerical values are INHERITED from specific microstructure or device-geometry choices made by the engineer.

\begin{itemize}[noitemsep]
\item Drude plasma frequency $\omega_p^2 = 2\pi c^2/[a^2\ln(a/r_0)]$ (\S4.4) --- Drude FORM FORCED; coefficient INHERITED from wire-array geometry.
\item Lorentz parameters $F$, $\omega_0$, $\gamma$ (\S4.5) --- Lorentz FORM FORCED; oscillator parameters INHERITED from SRR geometry.
\item Cloak parameters $\alpha$, $\beta$ (\S5.7) --- INHERITED from cloak-shell radii $R_1, R_2$.
\item Metasurface phase profile $\Phi(r_\parallel)$ (\S6.3) --- INHERITED from meta-atom geometry and arrangement.
\item Band structure $\omega_n(k)$ (\S7.3) --- FORM FORCED; specific band positions INHERITED from periodic geometry and material contrast.
\item Bandgap widths and locations (\S7.4) --- INHERITED from periodic geometry.
\item Corrector fields $\chi^j(y)$ (\S3.3) --- equation FORCED; specific solutions INHERITED from cell coefficient $A^{ij}(y)$.
\end{itemize}

\subsection{OPEN}

The following items are OPEN at the substrate level. They are not derived in this walkthrough and are flagged for future work.

\begin{itemize}[noitemsep]
\item Higher-order homogenization corrections at $a \sim \lambda/3$ to $\lambda$. The leading expansion is FORCED; next-order corrections at $(a/\lambda)^2$ are FORM-FORCED-VALUES-INHERITED from specific cell solutions but not derived here.
\item Nonlocal / spatially dispersive metamaterials. The closure machinery is local; spatial dispersion enters at the next order in $a/\lambda$ and requires extension beyond two-scale expansion.
\item Active and time-varying metamaterials. Time-dependent deformation $u(X, t)$ and time-varying microstructure are allowed by P-MM-1 and P-MM-3 respectively, but the corresponding closures (Floquet-in-time, parametric amplification, non-reciprocal phenomena) require their own derivations.
\item Topological photonics within Closure 1. Periodic substrates with engineered band topology, Berry curvature on the BZ, Chern numbers, photonic topological insulators --- adjacent to Closure 1 but not derived here.
\item Quantum-optical metamaterials. Single-photon and entangled-photon propagation through metamaterials uses the same coarse-grained wave equation at the field level; substrate-level quantum-state evolution through Closure 1/2/3 devices is reachable but not derived here.
\item Bulk-boundary correspondence for photonic systems. The relation between bulk band topology and boundary states is not derived here.
\item Hyperbolic metamaterials. Substrates with effective tensors of mixed-signature signature ($\varepsilon_\parallel > 0$, $\varepsilon_\perp < 0$ or vice versa) are within Closure 1's reach but not derived in detail here.
\item Time-domain metasurfaces. Spatiotemporal phase profiles $\Phi(r_\parallel, t)$ produce frequency shifts in addition to angle shifts; the static derivation of \S6 does not cover this case.
\end{itemize}

\textbf{No new substrate primitives are introduced anywhere in the walkthrough.}

\section{Exact Claims}

The substrate-level claims established in this walkthrough are:

\begin{enumerate}[noitemsep]
\item The substrate ontology FORCES a coarse-grained second-order wave equation (P-MM-4) whose coefficients are determined by the substrate's local rule-type content (P-MM-1).
\item Subwavelength periodic microstructure (P-MM-2) admits two-scale homogenization. The cell-problem corrector field $\chi^j(y)$, defined by an elliptic boundary-value problem on the unit cell with periodic boundary conditions and zero mean, produces effective constitutive tensors
\[
\varepsilon_{\rm eff}^{ij} = \langle\varepsilon^{ij}\rangle + \langle\varepsilon^{ik}\partial_{y^k}\chi_E^j\rangle, \qquad (\mu_{\rm eff}^{-1})^{ij} = \langle(\mu^{-1})^{ij}\rangle + \langle(\mu^{-1})^{ik}\partial_{y^k}\chi_M^j\rangle.
\]
\item The effective tensors satisfy Voigt-Reuss bounds. Microstructure produces anisotropy from isotropic constituents.
\item Engineered wire-array microstructures produce a Drude effective dielectric function with plasma frequency $\omega_p^2 = 2\pi c^2/[a^2\ln(a/r_0)]$, giving $\varepsilon_{\rm eff}(\omega) < 0$ for $\omega < \omega_p$.
\item Engineered split-ring resonator microstructures produce a Lorentz effective magnetic function with $\mu_{\rm eff}(\omega) < 0$ in a band immediately above the resonance frequency $\omega_0$.
\item In the doubly-negative band, the negative refractive index $n_{\rm eff} = -\sqrt{\varepsilon_r\mu_r}$ is FORCED by causal outgoing-wave selection. No separate postulate is required.
\item Slow substrate deformation (P-MM-3) produces transformed constitutive tensors $\varepsilon'^{ij} = (\det J)^{-1} J^i{}_k J^j{}_l \varepsilon^{kl}$ and an effective metric $g_{ij} = (J^{-1})^k{}_i (J^{-1})^l{}_j \delta_{kl}$ on the wave equation, written in covariant Laplace-Beltrami form.
\item The Pendry spherical cloak deformation $r' = R_1 + \alpha r$ with $\alpha = (R_2 - R_1)/R_2$ produces constitutive tensors $\varepsilon_{\hat r'}/\varepsilon_0 = \beta(r' - R_1)^2/r'^2$, $\varepsilon_{\hat\theta'}/\varepsilon_0 = \varepsilon_{\hat\phi'}/\varepsilon_0 = \beta$. The cloaked interior is topologically excluded at substrate level by the vanishing of $g_{\hat r'\hat r'}$ at $r' = R_1$.
\item The effective metric $g_{ij}$ is a coarse-grained rule-type metric on the substrate. It is \emph{not} spacetime curvature; no claim about gravitational physics is made or required.
\item Engineered codim-1 phase discontinuities (P-MM-5) imprint a tangential momentum kick $k_{\parallel,2} = k_{\parallel,1} + \nabla_\parallel\Phi$ on transmitted channels. The generalized Snell's law $n_1 \sin\theta_i = n_2 \sin\theta_t + (\lambda/2\pi)(d\Phi/dx)$ follows directly.
\item Periodic substrate at $a \sim \lambda$ supports Bloch eigenmodes $\psi_{n, k} = u_{n, k} e^{i k\cdot x}$ with band structure $\omega_n(k)$ on the Brillouin zone. Photonic bandgaps are frequency intervals where no eigenvalue exists, FORCED-FORM at substrate level.
\item The three substrate-level wave-control axes (microstructure, deformation, discontinuity) exhaust the precursor-machinery design space within P-MM-6.
\item Practical metamaterials are engineered compositions along the three axes. The unified substrate-level wave equation is
\[
\frac{1}{\sqrt{|g|}}\partial_i\!\left[\sqrt{|g|}\,g^{ij}(\varepsilon_{\rm eff}^{-1})^{jk}\partial_k\psi\right] + k_0^2\mu_{\rm eff}\psi = 0,
\]
with jump conditions $\psi(0^+) = \psi(0^-) e^{i\Phi(r_\parallel)}$ on engineered codim-1 surfaces.
\item The precursor machinery is bounded by the coarse-graining window P-MM-6, Kramers-Kronig dispersion-bandwidth trade-offs, and the Sommerfeld-Brillouin signal-front bound $v_{\rm front} \le c$.
\item The four Nobel-frontier metamaterials and photonics directions --- Yablonovitch's photonic bandgaps, Pendry's negative refractive index, Pendry's electromagnetic cloak, and Capasso's generalized Snell's law --- are reproduced at the substrate level by the three precursor closures plus the adjacent Bloch / bandgap regime, with no new substrate primitives.
\end{enumerate}

\section{References}

\begin{itemize}[noitemsep]
\item Yablonovitch, E. \textit{Inhibited spontaneous emission in solid-state physics and electronics.} Phys.\ Rev.\ Lett.\ \textbf{58}, 2059 (1987) --- photonic bandgaps.
\item John, S. \textit{Strong localization of photons in certain disordered dielectric superlattices.} Phys.\ Rev.\ Lett.\ \textbf{58}, 2486 (1987) --- companion bandgap paper.
\item Veselago, V. G. \textit{The electrodynamics of substances with simultaneously negative values of $\varepsilon$ and $\mu$.} Sov.\ Phys.\ Uspekhi \textbf{10}, 509 (1968) --- original theoretical analysis of negative index.
\item Pendry, J. B., Holden, A. J., Stewart, W. J., Youngs, I. \textit{Extremely low frequency plasmons in metallic mesostructures.} Phys.\ Rev.\ Lett.\ \textbf{76}, 4773 (1996) --- wire-array effective plasma frequency.
\item Pendry, J. B., Holden, A. J., Robbins, D. J., Stewart, W. J. \textit{Magnetism from conductors and enhanced nonlinear phenomena.} IEEE Trans.\ Microw.\ Theory Tech.\ \textbf{47}, 2075 (1999) --- split-ring resonator magnetic response.
\item Smith, D. R., Padilla, W. J., Vier, D. C., Nemat-Nasser, S. C., Schultz, S. \textit{Composite medium with simultaneously negative permeability and permittivity.} Phys.\ Rev.\ Lett.\ \textbf{84}, 4184 (2000) --- first experimental demonstration.
\item Pendry, J. B. \textit{Negative refraction makes a perfect lens.} Phys.\ Rev.\ Lett.\ \textbf{85}, 3966 (2000) --- negative refractive index and superlens.
\item Pendry, J. B., Schurig, D., Smith, D. R. \textit{Controlling electromagnetic fields.} Science \textbf{312}, 1780 (2006) --- transformation optics; spherical cloak.
\item Leonhardt, U. \textit{Optical conformal mapping.} Science \textbf{312}, 1777 (2006) --- concurrent transformation-optics construction.
\item Schurig, D., Mock, J. J., Justice, B. J., Cummer, S. A., Pendry, J. B., Starr, A. F., Smith, D. R. \textit{Metamaterial electromagnetic cloak at microwave frequencies.} Science \textbf{314}, 977 (2006) --- experimental realization.
\item Yu, N., Genevet, P., Kats, M. A., Aieta, F., Tetienne, J.-P., Capasso, F., Gaburro, Z. \textit{Light propagation with phase discontinuities: generalized laws of reflection and refraction.} Science \textbf{334}, 333 (2011) --- metasurfaces and generalized Snell's law.
\item Aieta, F., Genevet, P., Kats, M. A., Yu, N., Blanchard, R., Gaburro, Z., Capasso, F. \textit{Aberration-free ultrathin flat lenses and axicons at telecom wavelengths based on plasmonic metasurfaces.} Nano Lett.\ \textbf{12}, 4932 (2012) --- metasurface lens.
\item Bensoussan, A., Lions, J.-L., Papanicolaou, G. \textit{Asymptotic Analysis for Periodic Structures.} North-Holland (1978) --- mathematical foundations of two-scale homogenization.
\item Milton, G. W. \textit{The Theory of Composites.} Cambridge University Press (2002) --- homogenization, Voigt-Reuss bounds, effective tensors.
\item Joannopoulos, J. D., Johnson, S. G., Winn, J. N., Meade, R. D. \textit{Photonic Crystals: Molding the Flow of Light.} 2nd ed., Princeton University Press (2008) --- photonic-crystal band theory.
\item Sakoda, K. \textit{Optical Properties of Photonic Crystals.} 2nd ed., Springer (2005) --- Bloch analysis of photonic crystals.
\end{itemize}

\end{document}
